% 图 4.1 TVM 优化编译器框架转换步骤 — 基于 CICC0903540 技术文档 4.1
\documentclass[tikz,border=8pt]{standalone}
% 顶会/顶刊风格：低饱和度马卡龙 + 高级感
% 适用于 NeurIPS, CVPR, ICLR, IEEE TPAMI
\usepackage{amssymb}
\usepackage{fontspec}
\usepackage{xeCJK}
\setCJKmainfont{SimHei}
\usepackage{tikz}
\usepackage{pgfplots}
\usetikzlibrary{
  arrows.meta, positioning, shapes.geometric, calc,
  backgrounds, fit, decorations.pathreplacing, patterns
}

% 低饱和度马卡龙配色
\definecolor{mBlue}{HTML}{AEC6CF}
\definecolor{mPink}{HTML}{FFD1DC}
\definecolor{mGreen}{HTML}{B1D8B7}
\definecolor{mPurple}{HTML}{B39EB5}
\definecolor{mYellow}{HTML}{FDFD96}
\definecolor{mGray}{HTML}{E5E4E2}
\definecolor{mDark}{HTML}{4A4A4A}

% 全局 TikZ 样式
\tikzset{
  >=stealth,
  every node/.style={align=center, font=\sffamily},
  block/.style={
    rounded corners=3pt, draw=mDark, align=center,
    minimum height=0.8cm, inner sep=5pt},
  arr/.style={-{Stealth[length=4pt]}, semithick, draw=mDark},
  % 信息密度增强：张量维度标注
  ts/.style={font=\tiny\sffamily, color=mDark!80},
  % 数学操作符节点（⊕加 ⊗乘 ‖拼接）
  op/.style={circle, draw=mDark, minimum size=0.4cm, inner sep=0,
    font=\scriptsize\sffamily, fill=mGray!60},
  % 图例
  legendbox/.style={rectangle, draw=mDark!50, rounded corners=2pt,
    fill=mGray!15, inner sep=4pt, font=\tiny\sffamily},
}

\begin{document}
\begin{tikzpicture}[
  node distance=0.55cm and 0.6cm,
  stage/.style={
    rectangle, rounded corners=3pt, draw=mDark, align=center,
    font=\small\sffamily, semithick,
    minimum width=1.75cm, minimum height=0.7cm, inner sep=4pt},
  arr/.style={-{Stealth[length=4pt]}, draw=mDark, semithick},
  desc/.style={font=\scriptsize\sffamily, color=mDark, align=center, text width=1.7cm},
]
  \node[stage, fill=mBlue!40] (s1) {模型导入\\{\footnotesize TF/PyTorch/ONNX}};
  \node[stage, fill=mBlue!40, right=0.6cm of s1] (s2) {Relay/Relax\\{\footnotesize 高层次 IR}};
  \node[stage, fill=mBlue!40, right=0.6cm of s2] (s3) {TE\\{\footnotesize 张量表达式}};
  \node[stage, fill=mGreen!50, right=0.6cm of s3] (s4) {Auto-tuning\\{\footnotesize MetaSchedule}};
  \node[stage, fill=mBlue!40, right=0.6cm of s4] (s5) {TIR\\{\footnotesize 底层 IR}};
  \node[stage, fill=mPurple!40, right=0.6cm of s5] (s6) {机器码\\{\footnotesize .so}};

  \draw[arr] (s1) -- (s2)
    node[pos=0.5, above=2pt, ts] {计算图};
  \draw[arr] (s2) -- (s3)
    node[pos=0.5, above=2pt, ts] {FuseOps};
  \draw[arr] (s3) -- (s4)
    node[pos=0.5, above=2pt, ts] {schedule};
  \draw[arr] (s4) -- (s5)
    node[pos=0.5, above=2pt, ts] {JSON};
  \draw[arr] (s5) -- (s6)
    node[pos=0.5, above=2pt, ts] {LLVM};

  % 下方说明与张量标注分两行，拉开垂直间距
  \node[desc, below=0.35cm of s1] {从训练框架\\提取计算图};
  \node[desc, below=0.35cm of s2] {图级优化\\算子融合};
  \node[desc, below=0.35cm of s3] {循环切分\\矢量化};
  \node[desc, below=0.35cm of s4] {搜索最优\\schedule};
  \node[desc, below=0.35cm of s5] {底层优化\\内存重排};
  \node[desc, below=0.35cm of s6] {AArch64\\cortex-a72};

  \draw[decorate, decoration={brace, amplitude=3pt, raise=2pt}, draw=mDark, semithick]
    (s1.north west) -- (s6.north east)
    node[midway, above=6pt, font=\small\sffamily\bfseries, color=mDark]
    {TVM 端到端编译流水线};

  % 张量形状放在最下方一行，与 desc 有足够间距
  \node[ts, below=0.7cm of s1] {输入 $1{\times}32{\times}32{\times}32$};
  \node[ts, below=0.7cm of s6] {输出 $1{\times}3{\times}256{\times}256$};

  % 图例（右下，与内容留距）
  \node[legendbox, anchor=north east] at ([xshift=-4pt, yshift=-4pt]current bounding box.south east)
    [align=left, text width=3cm] {
    \textbf{图例}\\[1pt]
    \colorbox{mBlue!40}{\rule{0.14cm}{0.1cm}\hspace{0.05cm}} IR \quad
    \colorbox{mGreen!50}{\rule{0.14cm}{0.1cm}\hspace{0.05cm}} 调优 \quad
    \colorbox{mPurple!40}{\rule{0.14cm}{0.1cm}\hspace{0.05cm}} 产物
  };
\end{tikzpicture}
\end{document}
